\chapter*{Preface}
\addcontentsline{toc}{chapter}{Preface}

Most books about Vim are organized as inventories. They list commands, group them loosely by function, and trust the reader to assemble a working mental model from the accumulation of examples. This approach produces competent users of individual commands and, too often, users who never notice that the commands were never independent facts to begin with. They were instances of a small number of rules, applied over and over, in different combinations, against different kinds of text.

This book takes the opposite approach. It begins not with commands but with the entities Vim assumes exist: buffers, windows, registers, marks, modes, and the regions of text those modes and marks can address. It then derives the grammar that lets those entities combine, and only after the grammar is established does it turn to the specific transformations, that grammar makes possible. The intent is that a reader who finishes Part II should already be able to predict commands they have never seen, because they understand the rule that generates them rather than having memorized the rule's outputs.

The organizing progression is Ontology, Grammar, Geometry, Transformation, State, Automation, Application, and Reference. Each part depends on the one before it. Coordinate systems (Geometry) presuppose that the reader already knows what a motion is (Grammar), which presupposes knowing what a buffer is (Ontology). Automation presupposes a working vocabulary of transformations. Nothing in the main text is introduced before the concept it depends on has been established, with one deliberate exception: the encyclopedic appendices at the end, which are reference material rather than argument, and which the reader is free to consult out of order at any point.

This book is also deliberately conservative about scope. It covers core Vim: the editor as it has existed, largely unchanged in its fundamentals, for decades. Plugin ecosystems rise and fall; the operator-motion grammar, the register model, and the Ex language have not meaningfully changed since long before most current plugins existed. Neovim and its Lua configuration layer are mentioned where they diverge from core Vim in ways a reader should know about, but they are treated as an optional extension of the material here, not a parallel subject requiring separate treatment. A reader who learns core Vim from this book will be able to pick up Neovim, or the next editor built on the same ideas, without relearning the underlying grammar.

The practical workflows referenced throughout, particularly in the chapters on the Ex language and on Unix integration, are drawn from ordinary daily use rather than invented for demonstration: cleaning up manuscript text, editing configuration files, working inside Git, generating repetitive content, and the like. Vim is treated here not as an application with a feature list but as a text-processing language that happens to have a screen attached to it, one that composes naturally with the rest of the Unix toolchain rather than trying to replace it.

Anyone can read this book from the beginning and learn Vim from first principles. Anyone who already knows Vim reasonably well can instead read Part II and the early chapters of Part V, where the underlying grammar is made explicit, and will likely find that commands they already use unreflectively suddenly make more sense as instances of a rule rather than as isolated facts to be recalled individually.
